\documentclass[10pt,a4paper]{physicsquiz}
\tcbuselibrary{breakable,skins}
\setlength{\parindent}{0pt}
\setlength{\parskip}{0.35em}
\raggedcolumns%

\quizfontsize{9.4pt}{11.5pt}
\quiztitle{Comprehensive Practice: 60 Questions}%
\quizcourse{PHY 104  Oscillations, Normal Modes, Waves and Optics}
\quizauthor{Prepared by OT}
\quizinstitution{Drill Tutors}
\quizinstructions{Choose the best answer in each case. Questions 1—20 test foundations, Questions 21—40 test applied reasoning, and Questions 41—60 are multi-step challenges. Work independently before consulting the answer key and fully worked solutions.}
\quizconstants{\(g=\SI{9.80}{m.s^{-2}}\), \(v_{\text{sound}}=\SI{343}{m.s^{-1}}\) unless otherwise stated, \(c=\SI{3.00e8}{m.s^{-1}}\), \(\rho_{\text{air}}=\SI{1.21}{kg.m^{-3}}\), \(I_0=\SI{1.00e-12}{W.m^{-2}}\). Angles are measured from the normal unless stated otherwise.}

\newtcolorbox{workedsolution}[2]{
  breakable,
  enhanced,
  colback=white,
  colframe=QuizBlue,
  colbacktitle=QuizLightBlue,
  coltitle=QuizNavy,
  fonttitle=\bfseries,
  title={Question #1 — Correct option #2},
  arc=2mm,
  boxrule=0.65pt,
  left=8pt,right=8pt,top=6pt,bottom=6pt,
  before skip=8pt,after skip=8pt
}

\newcommand{\levelbanner}[2]{%
\item[]\begin{tcolorbox}[colback=QuizLightBlue,colframe=QuizBlue,boxrule=0.6pt,arc=1.5mm,left=5pt,right=5pt,top=4pt,bottom=4pt]
\centering\textbf{Level #1: #2}
\end{tcolorbox}}

\begin{document}
\pagenumbering{gobble}
\makequiztitle%
\pagenumbering{roman}
\tableofcontents
\clearpage
\pagenumbering{arabic}

\section{Question Booklet}
\constantsbox%
\begin{quizquestions}[2]
  \levelbanner{I}{Foundational Reasoning — Questions 1—20}
  \item A particle is displaced a distance \(x>0\) from equilibrium and executes simple harmonic motion. Which expression correctly gives the restoring force and its direction?
  \choices{\(F=+kx\), directed away from equilibrium}{\(F=-kx\), directed away from equilibrium}{\(F=-kx\), directed toward equilibrium}{\(F=-k/x\), directed toward equilibrium}{\(F=0\) because the particle is momentarily displaced}
  \item A \SI{0.50}{kg} mass is attached to a spring of constant \SI{200}{N.m^{-1}}. What is the period of the motion?
  \choices{\SI{0.157}{s}}{\SI{0.314}{s}}{\SI{0.628}{s}}{\SI{1.57}{s}}{\SI{3.14}{s}}
  \item A platform oscillates with amplitude \SI{4.0}{cm} and frequency \SI{2.5}{Hz}. What is its maximum speed?
  \choices{\SI{0.100}{m.s^{-1}}}{\SI{0.251}{m.s^{-1}}}{\SI{0.400}{m.s^{-1}}}{\SI{0.628}{m.s^{-1}}}{\SI{1.57}{m.s^{-1}}}
  \item A mass—spring oscillator has amplitude \(A\) and total mechanical energy \(E\). When the displacement is \(x=A/2\), which statement is correct?
  \choices{\(U=E/2\) and \(K=E/2\)}{\(U=E/4\) and \(K=3E/4\)}{\(U=3E/4\) and \(K=E/4\)}{\(U=E\) and \(K=0\)}{\(U=0\) and \(K=E\)}
  \item What length must a simple pendulum have to possess a period of \SI{2.00}{s} near Earth's surface?
  \choices{\SI{0.993}{m}}{\SI{1.56}{m}}{\SI{2.00}{m}}{\SI{3.12}{m}}{\SI{9.80}{m}}
  \item Which statement best defines a normal mode of a coupled system?
  \choices{Every part remains at rest while energy moves through the system.}{Every part oscillates with one angular frequency and fixed phase relations.}{Different parts oscillate at unrelated frequencies.}{Only one part of the system is allowed to move.}{The coupling force must vanish at every instant.}
  \item Two identical masses \(m\) are attached to end springs of constant \(k\) and joined by a coupling spring \(k_c\). What is the angular frequency of the in-phase mode?
  \choices{\(\sqrt{k/m}\)}{\(\sqrt{k_c/m}\)}{\(\sqrt{(k+k_c)/m}\)}{\(\sqrt{(k+2k_c)/m}\)}{\(\sqrt{2k/m}\)}
  \item Which amplitude vector represents the out-of-phase normal mode of two identical coupled oscillators?
  \choices{\(\begin{bmatrix}0\\0\end{bmatrix}\)}{\(\begin{bmatrix}1\\0\end{bmatrix}\)}{\(\begin{bmatrix}1\\1\end{bmatrix}\)}{\(\begin{bmatrix}1\\-1\end{bmatrix}\)}{\(\begin{bmatrix}1\\2\end{bmatrix}\)}
  \item Why is the out-of-phase normal frequency higher than the in-phase normal frequency in the symmetric two-mass system?
  \choices{The masses become heavier in the out-of-phase mode.}{The end springs stop exerting forces.}{The coupling spring is strongly stretched or compressed and adds restoring force.}{The coupling spring removes kinetic energy permanently.}{The amplitude vector is not an eigenvector.}
  \item For \(k=\SI{100}{N.m^{-1}}\) and \(k_c=\SI{50}{N.m^{-1}}\), what is the ratio \(\omega_2/\omega_1\)?
  \choices{\(1.00\)}{\(\sqrt{2}\approx1.41\)}{\(2.00\)}{\(\sqrt{3}\approx1.73\)}{\(0.707\)}
  \item A transverse wave is described by \(y=(\SI{3.0}{mm})\sin(4x-20t)\), with SI units. Which pair gives its wavelength and speed?
  \choices{\(\lambda=\SI{0.250}{m}\), \(v=\SI{80}{m.s^{-1}}\)}{\(\lambda=\SI{1.00}{m}\), \(v=\SI{5.0}{m.s^{-1}}\)}{\(\lambda=\SI{1.57}{m}\), \(v=\SI{5.0}{m.s^{-1}}\)}{\(\lambda=\SI{1.57}{m}\), \(v=\SI{80}{m.s^{-1}}\)}{\(\lambda=\SI{6.28}{m}\), \(v=\SI{0.20}{m.s^{-1}}\)}
  \item The wave \(y=A\cos(6x+24t)\) travels in which direction and at what speed?
  \choices{+\(x\) at \SI{24}{m.s^{-1}}}{+\(x\) at \SI{4}{m.s^{-1}}}{-\(x\) at \SI{24}{m.s^{-1}}}{-\(x\) at \SI{4}{m.s^{-1}}}{It is a standing wave.}
  \item A standing wave has wavelength \SI{0.80}{m}. What is the separation between adjacent nodes?
  \choices{\SI{0.20}{m}}{\SI{0.40}{m}}{\SI{0.80}{m}}{\SI{1.20}{m}}{\SI{1.60}{m}}
  \item A string of length \SI{0.80}{m}, fixed at both ends, carries waves at \SI{120}{m.s^{-1}}. What is the third-harmonic frequency?
  \choices{\SI{75}{Hz}}{\SI{120}{Hz}}{\SI{150}{Hz}}{\SI{180}{Hz}}{\SI{225}{Hz}}
  \item A sound intensity is \SI{1.0e-6}{W.m^{-2}}. What is the corresponding sound level? Take \(I_0=\SI{1.0e-12}{W.m^{-2}}\).
  \choices{\SI{6}{dB}}{\SI{12}{dB}}{\SI{60}{dB}}{\SI{100}{dB}}{\SI{120}{dB}}
  \item Light of vacuum wavelength \SI{600}{nm} enters a medium of refractive index \(1.50\). Which pair gives its speed and wavelength in the medium?
  \choices{\(\SI{2.00e8}{m.s^{-1}}\) and \SI{400}{nm}}{\(\SI{2.00e8}{m.s^{-1}}\) and \SI{600}{nm}}{\(\SI{3.00e8}{m.s^{-1}}\) and \SI{400}{nm}}{\(\SI{4.50e8}{m.s^{-1}}\) and \SI{900}{nm}}{\(\SI{1.50e8}{m.s^{-1}}\) and \SI{300}{nm}}
  \item A ray passes from air into glass of refractive index \(1.50\) at an incidence angle of \(30.0^\circ\). What is the refracted angle?
  \choices{\(10.0^\circ\)}{\(15.0^\circ\)}{\(30.0^\circ\)}{\(19.5^\circ\)}{\(48.6^\circ\)}
  \item In a Young double-slit experiment, \(\lambda=\SI{600}{nm}\), the screen distance is \SI{2.0}{m}, and the slit separation is \SI{0.40}{mm}. What is the fringe spacing?
  \choices{\SI{0.75}{mm}}{\SI{3.00}{mm}}{\SI{6.00}{mm}}{\SI{12.0}{mm}}{\SI{30.0}{mm}}
  \item Light of wavelength \SI{500}{nm} passes through a single slit of width \SI{0.20}{mm}. At what angle does the first minimum occur?
  \choices{\(0.0143^\circ\)}{\(0.0286^\circ\)}{\(0.0716^\circ\)}{\(0.100^\circ\)}{\(0.143^\circ\)}
  \item A telescope of aperture diameter \SI{0.10}{m} observes light of wavelength \SI{550}{nm}. What is its diffraction-limited angular resolution?
  \choices{\(5.50\times10^{-7}\ \mathrm{rad}\)}{\(4.51\times10^{-6}\ \mathrm{rad}\)}{\(6.71\times10^{-6}\ \mathrm{rad}\)}{\(1.22\times10^{-5}\ \mathrm{rad}\)}{\(6.71\times10^{-5}\ \mathrm{rad}\)}
  \levelbanner{II}{Applied Problems — Questions 21—40}
  \item An oscillator obeys \(x=A\cos(\omega t+\phi)\). At \(t=0\), \(x=\SI{3.0}{cm}\) and \(v=\SI{-0.40}{m.s^{-1}}\), while \(\omega=\SI{10}{rad.s^{-1}}\). What are \(A\) and \(\phi\)?
  \choices{\(A=\SI{4.0}{cm}\), \(\phi=36.9^\circ\)}{\(A=\SI{5.0}{cm}\), \(\phi=53.1^\circ\)}{\(A=\SI{5.0}{cm}\), \(\phi=-53.1^\circ\)}{\(A=\SI{7.0}{cm}\), \(\phi=45.0^\circ\)}{\(A=\SI{3.0}{cm}\), \(\phi=90.0^\circ\)}
  \item A uniform rod of length \SI{1.20}{m} swings with small amplitude about a horizontal pivot through one end. What is its period?
  \choices{\SI{0.898}{s}}{\SI{1.27}{s}}{\SI{1.55}{s}}{\SI{1.80}{s}}{\SI{2.20}{s}}
  \item Three identical springs, each of constant \SI{120}{N.m^{-1}}, are connected in series to a \SI{0.36}{kg} mass. What is the oscillation period?
  \choices{\SI{0.172}{s}}{\SI{0.344}{s}}{\SI{0.596}{s}}{\SI{1.03}{s}}{\SI{1.88}{s}}
  \item A \SI{10.0}{g} bullet moving at \SI{400}{m.s^{-1}} embeds in a \SI{1.99}{kg} block attached to a spring of constant \SI{800}{N.m^{-1}}. The block is initially at equilibrium and at rest. Which pair gives the speed immediately after impact and the resulting amplitude?
  \choices{\SI{2.00}{m.s^{-1}} and \SI{0.100}{m}}{\SI{4.00}{m.s^{-1}} and \SI{0.200}{m}}{\SI{2.01}{m.s^{-1}} and \SI{0.050}{m}}{\SI{1.00}{m.s^{-1}} and \SI{0.100}{m}}{\SI{200}{m.s^{-1}} and \SI{0.500}{m}}
  \item A torsional oscillator has torsion constant \[\kappa=\SI{0.80}{N.m.rad^{-1}},\] moment of inertia \(I=\SI{0.020}{kg.m^2}\), and angular amplitude \(\theta_{\max}=\SI{0.25}{rad}\). What is its angular speed when \(\theta=\SI{0.15}{rad}\)?
  \choices{\SI{0.400}{rad.s^{-1}}}{\SI{0.800}{rad.s^{-1}}}{\SI{1.00}{rad.s^{-1}}}{\SI{1.60}{rad.s^{-1}}}{\SI{1.26}{rad.s^{-1}}}
  \item For two identical coupled oscillators, \(m=\SI{0.50}{kg}\), \(k=\SI{100}{N.m^{-1}}\), and \(k_c=\SI{50}{N.m^{-1}}\). What are the two normal angular frequencies?
  \choices{\(10.0\) and \(14.1\ \mathrm{rad\,s^{-1}}\)}{\(14.1\) and \(20.0\ \mathrm{rad\,s^{-1}}\)}{\(20.0\) and \(28.3\ \mathrm{rad\,s^{-1}}\)}{\(7.07\) and \(14.1\ \mathrm{rad\,s^{-1}}\)}{\(14.1\) and \(17.3\ \mathrm{rad\,s^{-1}}\)}
  \item A symmetric two-mass system has \(m=\SI{2.0}{kg}\) and measured normal angular frequencies \(\omega_1=\SI{10}{rad.s^{-1}}\) and \(\omega_2=\SI{14}{rad.s^{-1}}\). Find \(k\) and \(k_c\).
  \choices{\(k=\SI{100}{N.m^{-1}}\), \(k_c=\SI{48}{N.m^{-1}}\)}{\(k=\SI{196}{N.m^{-1}}\), \(k_c=\SI{100}{N.m^{-1}}\)}{\(k=\SI{200}{N.m^{-1}}\), \(k_c=\SI{192}{N.m^{-1}}\)}{\(k=\SI{200}{N.m^{-1}}\), \(k_c=\SI{96}{N.m^{-1}}\)}{\(k=\SI{392}{N.m^{-1}}\), \(k_c=\SI{96}{N.m^{-1}}\)}
  \item The initial displacement of two identical coupled masses is \(\mathbf{x}(0)=\begin{bmatrix}A\\0\end{bmatrix}\), with zero initial velocities. In the expansion \(\mathbf{x}=C_1\begin{bmatrix}1\\1\end{bmatrix}+C_2\begin{bmatrix}1\\-1\end{bmatrix}\), what are \(C_1\) and \(C_2\) at \(t=0\)?
  \choices{\(C_1=A\), \(C_2=0\)}{\(C_1=0\), \(C_2=A\)}{\(C_1=A/2\), \(C_2=A/2\)}{\(C_1=A\), \(C_2=-A\)}{\(C_1=2A\), \(C_2=2A\)}
  \item For \(x_1=\SI{3.0}{cm}\) and \(x_2=\SI{-1.0}{cm}\), what are the normal coordinates \(q_1=x_1+x_2\) and \(q_2=x_1-x_2\)?
  \choices{\(q_1=\SI{2.0}{cm}\), \(q_2=\SI{4.0}{cm}\)}{\(q_1=\SI{4.0}{cm}\), \(q_2=\SI{2.0}{cm}\)}{\(q_1=\SI{2.0}{cm}\), \(q_2=\SI{-4.0}{cm}\)}{\(q_1=\SI{-2.0}{cm}\), \(q_2=\SI{4.0}{cm}\)}{\(q_1=\SI{3.0}{cm}\), \(q_2=\SI{-1.0}{cm}\)}
  \item Which characteristic equation determines the normal angular frequencies of the symmetric two-mass system?
  \choices{\(k+k_c-m\omega^2=0\)}{\({(k-m\omega^2)}^2+k_c^2=0\)}{\({(k+k_c+m\omega^2)}^2-k_c^2=0\)}{\({(k+2k_c-m\omega)}^2=0\)}{\({(k+k_c-m\omega^2)}^2-k_c^2=0\)}
  \item A string wave is \(y=(\SI{4.0}{mm})\sin(5x-100t)\). If the linear density is \SI{0.010}{kg.m^{-1}}, which set gives the wave speed, tension, and average power?
  \choices{\(\SI{5}{m.s^{-1}}\), \SI{0.25}{N}, \SI{0.004}{W}}{\(\SI{20}{m.s^{-1}}\), \SI{4.0}{N}, \SI{0.016}{W}}{\(\SI{20}{m.s^{-1}}\), \SI{2.0}{N}, \SI{0.032}{W}}{\(\SI{100}{m.s^{-1}}\), \SI{100}{N}, \SI{0.080}{W}}{\(\SI{500}{m.s^{-1}}\), \SI{2500}{N}, \SI{1.60}{W}}
  \item A pipe of length \SI{0.85}{m} is closed at one end and open at the other. Taking the sound speed as \SI{340}{m.s^{-1}}, what are its first three resonant frequencies?
  \choices{\(100,200,300\ \mathrm{Hz}\)}{\(200,400,600\ \mathrm{Hz}\)}{\(100,300,500\ \mathrm{Hz}\)}{\(50,150,250\ \mathrm{Hz}\)}{\(100,400,700\ \mathrm{Hz}\)}
  \item A stretched wire has fundamental frequency \SI{500}{Hz}. Its tension is increased slightly so that it produces \SI{5.0}{beats.s^{-1}} with the original wire. What fractional increase in tension is required?
  \choices{\(0.0050\)}{\(0.0100\)}{\(0.0150\)}{\(0.0201\)}{\(0.0404\)}
  \item A source emitting \SI{600}{Hz} moves directly toward a stationary observer at \SI{30}{m.s^{-1}}. Taking \(v=\SI{343}{m.s^{-1}}\), what frequency is observed?
  \choices{\SI{658}{Hz}}{\SI{600}{Hz}}{\SI{552}{Hz}}{\SI{630}{Hz}}{\SI{714}{Hz}}
  \item An aircraft travels at Mach \(1.50\). What is the Mach-cone half-angle?
  \choices{\(30.0^\circ\)}{\(33.7^\circ\)}{\(45.0^\circ\)}{\(60.0^\circ\)}{\(41.8^\circ\)}
  \item A mica sheet of refractive index \(1.60\) and thickness \SI{5.0}{\micro\,m} is placed over one slit in a double-slit experiment using \SI{600}{nm} light. By how many fringe spacings does the pattern shift?
  \choices{\(3\)}{\(5\)}{\(6\)}{\(8\)}{\(10\)}
  \item In a real double-slit arrangement, the slit separation is four times the slit width: \(d=4a\). Which interference orders are missing?
  \choices{\(m=1,2,3,\ldots\)}{\(m=2,4,6,\ldots\)}{\(m=4,8,12,\ldots\)}{\(m=3,6,9,\ldots\)}{No orders are missing.}
  \item A diffraction grating has \(500\) lines per millimetre and is illuminated with \SI{600}{nm} light. Which pair gives the second-order angle and the maximum possible order?
  \choices{\(17.5^\circ\) and \(m_{\max}=2\)}{\(30.0^\circ\) and \(m_{\max}=3\)}{\(53.1^\circ\) and \(m_{\max}=4\)}{\(36.9^\circ\) and \(m_{\max}=3\)}{\(90.0^\circ\) and \(m_{\max}=2\)}
  \item For glass of refractive index \(1.50\) surrounded by air, which pair gives (i) the critical angle for glass-to-air incidence and (ii) Brewster's angle for air-to-glass incidence?
  \choices{\(41.8^\circ\) and \(56.3^\circ\)}{\(56.3^\circ\) and \(41.8^\circ\)}{\(33.7^\circ\) and \(48.2^\circ\)}{\(48.2^\circ\) and \(33.7^\circ\)}{\(90.0^\circ\) and \(45.0^\circ\)}
  \item X-rays of wavelength \SI{0.154}{nm} produce a first-order Bragg peak at a measured scattering angle \(2\theta=50.0^\circ\). What is the crystal-plane spacing?
  \choices{\SI{0.0770}{nm}}{\SI{0.119}{nm}}{\SI{0.154}{nm}}{\SI{0.238}{nm}}{\SI{0.182}{nm}}
  \levelbanner{III}{Multi-step Challenges — Questions 41—60}
  \item An oscillator has amplitude \SI{6.0}{cm} and angular frequency \SI{8.0}{rad.s^{-1}}. At \(t=0\), its displacement is \SI{3.0}{cm} and it is moving in the positive direction. Which equation is correct?
  \choices{\(x=0.060\cos(8t+\pi/3)\)}{\(x=0.030\cos(8t-\pi/3)\)}{\(x=0.060\cos(8t-\pi/3)\)}{\(x=0.060\sin(8t-\pi/3)\)}{\(x=0.060\cos(8t+2\pi/3)\)}
  \item At a displacement of \SI{3.0}{cm}, a mass in SHM has speed \SI{0.40}{m.s^{-1}}. Its turning points are at \(x=\pm\SI{5.0}{cm}\). What is the period?
  \choices{\SI{0.314}{s}}{\SI{0.628}{s}}{\SI{0.785}{s}}{\SI{1.26}{s}}{\SI{1.57}{s}}
  \item A uniform rod of length \SI{1.20}{m} oscillates as a physical pendulum about a pivot located one-quarter of its length from one end. What is its small-angle period?
  \choices{\SI{0.840}{s}}{\SI{1.12}{s}}{\SI{1.44}{s}}{\SI{1.68}{s}}{\SI{2.06}{s}}
  \item At \(t=0\), an oscillator is at \(x=A/2\) and moving toward equilibrium. How long does it first take to reach the opposite turning point \(x=-A\)?
  \choices{\(T/3\)}{\(T/4\)}{\(T/2\)}{\(2T/3\)}{\(3T/4\)}
  \item A \SI{1.00}{kg} block attached to a \SI{100}{N.m^{-1}} spring passes through equilibrium at \SI{2.00}{m.s^{-1}}. At that instant, a \SI{0.25}{kg} lump of clay moving vertically falls onto the block and sticks. What is the new horizontal oscillation amplitude?
  \choices{\SI{0.100}{m}}{\SI{0.125}{m}}{\SI{0.160}{m}}{\SI{0.200}{m}}{\SI{0.179}{m}}
  \item For a symmetric coupled system with \(m=\SI{1}{kg}\), \(k=\SI{4}{N.m^{-1}}\), and \(k_c=\SI{6}{N.m^{-1}}\), mass 1 is initially displaced by \(A\), mass 2 is at equilibrium, and both are released from rest. What are the displacements at \(t=\pi/2\ \mathrm{s}\)?
  \choices{\(x_1=A\), \(x_2=0\)}{\(x_1=-A\), \(x_2=0\)}{\(x_1=0\), \(x_2=-A\)}{\(x_1=0\), \(x_2=A\)}{\(x_1=A/2\), \(x_2=-A/2\)}
  \item Two identical masses of \SI{0.50}{kg} have normal frequencies \(f_1=\SI{2.00}{Hz}\) and \(f_2=\SI{3.00}{Hz}\). What are \(k\) and \(k_c\)?
  \choices{\(k=\SI{19.7}{N.m^{-1}}\), \(k_c=\SI{12.3}{N.m^{-1}}\)}{\(k=\SI{79.0}{N.m^{-1}}\), \(k_c=\SI{49.3}{N.m^{-1}}\)}{\(k=\SI{78.9}{N.m^{-1}}\), \(k_c=\SI{98.7}{N.m^{-1}}\)}{\(k=\SI{157.9}{N.m^{-1}}\), \(k_c=\SI{49.3}{N.m^{-1}}\)}{\(k=\SI{39.5}{N.m^{-1}}\), \(k_c=\SI{24.7}{N.m^{-1}}\)}
  \item One mass of a weakly coupled identical pair is displaced and both are released from rest. If the normal angular frequencies are \SI{10}{rad.s^{-1}} and \SI{12}{rad.s^{-1}}, when does the first complete transfer of oscillation amplitude to the other mass occur?
  \choices{\SI{0.262}{s}}{\SI{0.524}{s}}{\SI{1.05}{s}}{\SI{1.57}{s}}{\SI{3.14}{s}}
  \item For the same displacement amplitude \(A\) of each mass, compare the total energies of the pure out-of-phase and in-phase modes when \(k=\SI{100}{N.m^{-1}}\) and \(k_c=\SI{50}{N.m^{-1}}\). What is \(E_{\mathrm{out}}/E_{\mathrm{in}}\)?
  \choices{\(2.00\)}{\(1.50\)}{\(1.00\)}{\(0.50\)}{\(3.00\)}
  \item A system with unit masses has stiffness matrix \(K=\begin{bmatrix}6&-2\\-2&6\end{bmatrix}\ \mathrm{N.m^{-1}}\). Its initial displacement is \(\begin{bmatrix}3\\1\end{bmatrix}\ \mathrm{cm}\). Which option gives the modal coefficients in the basis \({[1,1]}^T,{[1,-1]}^T\) and the corresponding angular frequencies?
  \choices{\(C_1=1\ \mathrm{cm}\), \(C_2=2\ \mathrm{cm}\); \(\omega_1=\sqrt6\), \(\omega_2=\sqrt2\)}{\(C_1=3\ \mathrm{cm}\), \(C_2=1\ \mathrm{cm}\); \(\omega_1=2\), \(\omega_2=4\)}{\(C_1=2\ \mathrm{cm}\), \(C_2=-1\ \mathrm{cm}\); \(\omega_1=\sqrt8\), \(\omega_2=2\)}{\(C_1=4\ \mathrm{cm}\), \(C_2=2\ \mathrm{cm}\); \(\omega_1=4\), \(\omega_2=8\)}{\(C_1=2\ \mathrm{cm}\), \(C_2=1\ \mathrm{cm}\); \(\omega_1=2\), \(\omega_2=2\sqrt2\)}
  \item Three waves of equal frequency and amplitude \SI{5.0}{mm} have phase constants \(0\), \(\pi/3\), and \(2\pi/3\). What are the amplitude and phase constant of their resultant?
  \choices{\SI{5.0}{mm} and \(0\)}{\SI{8.66}{mm} and \(\pi/2\)}{\SI{10.0}{mm} and \(\pi/3\)}{\SI{15.0}{mm} and \(\pi/3\)}{\(0\) and an undefined phase}
  \item A stone is dropped from rest into a well. The splash is heard \SI{4.00}{s} later. Taking \(g=\SI{9.80}{m.s^{-2}}\) and the sound speed as \SI{343}{m.s^{-1}}, what is the well depth?
  \choices{\SI{61.0}{m}}{\SI{70.5}{m}}{\SI{78.4}{m}}{\SI{86.0}{m}}{\SI{39.2}{m}}
  \item A stationary detector emits ultrasound at \SI{40.0}{kHz} toward a car approaching at \SI{30.0}{m.s^{-1}}. The wave reflects from the car and returns to the detector. Taking \(v=\SI{343}{m.s^{-1}}\), what return frequency is detected?
  \choices{\SI{36.8}{kHz}}{\SI{40.0}{kHz}}{\SI{43.5}{kHz}}{\SI{45.2}{kHz}}{\SI{47.7}{kHz}}
  \item A standing wave is \(y=0.080\sin(\pi x/2)\cos(4\pi t)\) in SI units. At \(x=\SI{0.50}{m}\) and \(t=\SI{0.0625}{s}\), which pair gives the local amplitude and instantaneous displacement?
  \choices{\SI{56.6}{mm} and \SI{40.0}{mm}}{\SI{80.0}{mm} and \SI{56.6}{mm}}{\SI{40.0}{mm} and \SI{56.6}{mm}}{\SI{56.6}{mm} and \(0\)}{\SI{40.0}{mm} and \SI{40.0}{mm}}
  \item Two coherent in-phase point sources are separated by \SI{1.50}{m} and emit wavelength \SI{0.400}{m}. A detector moves once around a very large circle in the plane of the sources. How many intensity maxima does it encounter?
  \choices{\(7\)}{\(8\)}{\(12\)}{\(14\)}{\(16\)}
  \item For a real double slit, the intensity is \[I=I_0{(\sin\beta/\beta)}^2\cos^2\alpha.\] At an angle where \(\beta=\pi/2\) and \(\alpha=\pi/3\), what is \(I/I_0\)?
  \choices{\(1/4\)}{\(1/\pi^2\approx0.101\)}{\(4/\pi^2\approx0.405\)}{\(1/2\)}{\(0\)}
  \item A grating must just resolve \SI{600.0}{nm} and \SI{600.3}{nm} in second order. If the grating has \(500\) lines per millimetre, what minimum number of illuminated lines and illuminated width are required?
  \choices{\(N=500\) and \SI{1.00}{mm}}{\(N=667\) and \SI{1.33}{mm}}{\(N=1000\) and \SI{2.00}{mm}}{\(N=2000\) and \SI{4.00}{mm}}{\(N=4000\) and \SI{8.00}{mm}}
  \item For a single slit of width \SI{0.20}{mm} illuminated by \SI{600}{nm} light, the first secondary maximum occurs at \(\beta\approx4.493\). Approximately where is it, and what fraction of the central intensity does it have?
  \choices{\(0.086^\circ\) and \(50\%\)}{\(0.143^\circ\) and \(10\%\)}{\(0.200^\circ\) and \(1.0\%\)}{\(0.246^\circ\) and \(4.7\%\)}{\(0.492^\circ\) and \(0.47\%\)}
  \item Two headlights separated by \SI{1.50}{m} are \SI{100}{km} away. What minimum telescope aperture is required to just resolve them at \SI{550}{nm}, according to Rayleigh's criterion?
  \choices{\SI{4.47}{cm}}{\SI{2.24}{cm}}{\SI{8.94}{cm}}{\SI{44.7}{cm}}{\SI{1.22}{cm}}
  \item A grating has \(300\) lines per millimetre, and each slit has width \(a=d/3\), where \(d\) is the grating spacing. It is illuminated with \SI{500}{nm} light. How many observable principal maxima occur over both sides, including the central maximum, after missing orders are excluded?
  \choices{\(5\)}{\(7\)}{\(11\)}{\(13\)}{\(9\)}
\end{quizquestions}
\clearpage

\section{Answer Key and Topic Map}
\begin{answerkey}
  \centering
  \renewcommand{\arraystretch}{1.35}
  \begin{tabular}{*{5}{>{\centering\arraybackslash}p{0.16\linewidth}}}
    \toprule
    \textbf{1.} C  & \textbf{2.} B  & \textbf{3.} D  & \textbf{4.} B  & \textbf{5.} A  \\
    \textbf{6.} B  & \textbf{7.} A  & \textbf{8.} D  & \textbf{9.} C  & \textbf{10.} B \\
    \textbf{11.} C & \textbf{12.} D & \textbf{13.} B & \textbf{14.} E & \textbf{15.} C \\
    \textbf{16.} A & \textbf{17.} D & \textbf{18.} B & \textbf{19.} E & \textbf{20.} C \\
    \textbf{21.} B & \textbf{22.} D & \textbf{23.} C & \textbf{24.} A & \textbf{25.} E \\
    \textbf{26.} B & \textbf{27.} D & \textbf{28.} C & \textbf{29.} A & \textbf{30.} E \\
    \textbf{31.} B & \textbf{32.} C & \textbf{33.} D & \textbf{34.} A & \textbf{35.} E \\
    \textbf{36.} B & \textbf{37.} C & \textbf{38.} D & \textbf{39.} A & \textbf{40.} E \\
    \textbf{41.} C & \textbf{42.} B & \textbf{43.} D & \textbf{44.} A & \textbf{45.} E \\
    \textbf{46.} C & \textbf{47.} B & \textbf{48.} D & \textbf{49.} A & \textbf{50.} E \\
    \textbf{51.} C & \textbf{52.} B & \textbf{53.} E & \textbf{54.} A & \textbf{55.} D \\
    \textbf{56.} B & \textbf{57.} C & \textbf{58.} D & \textbf{59.} A & \textbf{60.} E \\
    \bottomrule
  \end{tabular}
\end{answerkey}

\begin{tcolorbox}[colback=QuizLightGrey,colframe=QuizGrey,arc=2mm,boxrule=0.6pt]
  \textbf{Topic distribution:} Each topic contributes 15 questions: Oscillations (1--5, 21--25, 41--45), Normal Modes (6--10, 26--30, 46--50), Waves and Sound (11--15, 31--35, 51--55), and Optics (16--20, 36--40, 56--60).
\end{tcolorbox}

\clearpage
\section{Fully Worked Solutions}
\subsection{Level I:\ Foundational Reasoning}
\begin{workedsolution}{1}{C}\textbf{Topic: Oscillations.}\par For simple harmonic motion, the defining force law is Hooke's law,
  \[
    F=-kx.
  \]
  The magnitude \(kx\) is proportional to the displacement. The minus sign does not mean that the force is always numerically negative; it means that the force points opposite to the displacement. Since \(x>0\), the restoring force points in the negative direction, toward equilibrium. Hence option \textbf{C}.\end{workedsolution}
\begin{workedsolution}{2}{B}\textbf{Topic: Oscillations.}\par For a mass{--}spring oscillator,
  \[
    T=2\pi\sqrt{\frac{m}{k}}.
  \]
  Substituting \(m=\SI{0.50}{kg}\) and \(k=\SI{200}{N.m^{-1}}\),
  \[
    T=2\pi\sqrt{\frac{0.50}{200}}
    =2\pi\sqrt{0.0025}
    =2\pi(0.0500)
    =\SI{0.314}{s}.
  \]
  Therefore, option \textbf{B} is correct.\end{workedsolution}
\begin{workedsolution}{3}{D}\textbf{Topic: Oscillations.}\par The maximum speed in SHM is
  \[
    v_{\max}=\omega A=2\pi fA.
  \]
  Here, \(A=\SI{0.040}{m}\) and \(f=\SI{2.5}{Hz}\), so
  \[
    v_{\max}=2\pi(2.5)(0.040)=\SI{0.628}{m.s^{-1}}.
  \]
  Thus, option \textbf{D}.\end{workedsolution}
\begin{workedsolution}{4}{B}\textbf{Topic: Oscillations.}\par For a spring oscillator,
  \[
    E=\frac12 kA^2,\qquad U=\frac12 kx^2.
  \]
  At \(x=A/2\),
  \[
    U=\frac12 k{\left\{\frac{A}{2}\right\}}^2
    =\frac18 kA^2
    =\frac14\left(\frac12 kA^2\right)
    =\frac{E}{4}.
  \]
  Energy conservation gives \(K=E-U=3E/4\). Hence option \textbf{B}.\end{workedsolution}
\begin{workedsolution}{5}{A}\textbf{Topic: Oscillations.}\par For a small-amplitude simple pendulum,
  \[
    T=2\pi\sqrt{\frac{L}{g}}.
  \]
  Solving for \(L\),
  \[
    L=g{\left(\frac{T}{2\pi}\right)}^2
    =9.80{\left(\frac{2.00}{2\pi}\right)}^2
    =\SI{0.993}{m}.
  \]
  Therefore, option \textbf{A}.\end{workedsolution}
\begin{workedsolution}{6}{B}\textbf{Topic: Normal Modes.}\par A normal mode is a special coordinated motion. Every moving part has the same angular frequency, while the phase relation and amplitude ratio between parts remain fixed. The parts don't necessarily have equal amplitudes, and the coupling force need not always vanish. Hence option \textbf{B}.\end{workedsolution}
\begin{workedsolution}{7}{A}\textbf{Topic: Normal Modes.}\par In the in-phase mode, \(x_1=x_2\). The coupling spring has no change in length because both masses move together (in the same direction). It therefore adds no restoring force. Each mass behaves as though attached only to its outer spring, giving
  \[
    \omega_1=\sqrt{\frac{k}{m}}.
  \]
  Thus, option \textbf{A}.\end{workedsolution}
\begin{workedsolution}{8}{D}\textbf{Topic: Normal Modes.}\par Out-of-phase motion requires equal and opposite displacements:
  \[
    x_1=-x_2.
  \]
  The vector \({[1,-1]}^T\) encodes precisely this relation. Any nonzero scalar multiple, such as \({[3,-3]}^T\), describes the same mode shape.\\
  I may write vectors like \({[1, 1]}^T\), instead of the usual \({\displaystyle{\binom{1}{1}}}\) Therefore, option \textbf{D}.\end{workedsolution}
\begin{workedsolution}{9}{C}\textbf{Topic: Normal Modes.}\par When the masses move oppositely, their relative displacement is large. The middle spring is therefore stretched or compressed substantially, so it contributes an additional restoring force. A stronger effective restoring action produces a larger angular frequency. Hence option \textbf{C}.\end{workedsolution}
\begin{workedsolution}{10}{B}\textbf{Topic: Normal Modes.}\par The two angular frequencies are
  \[
    \omega_1=\sqrt{\frac{k}{m}},\qquad
    \omega_2=\sqrt{\frac{k+2k_c}{m}}.
  \]
  Thus,
  \[
    \frac{\omega_2}{\omega_1}
    =\sqrt{\frac{k+2k_c}{k}}
    =\sqrt{\frac{100+2(50)}{100}}
    =\sqrt{2}.
  \]
  Therefore, option \textbf{B}.\end{workedsolution}
\begin{workedsolution}{11}{C}\textbf{Topic: Waves.}\par Comparing with \(y=A\sin(kx-\omega t)\) gives \(k=\SI{4}{rad.m^{-1}}\) and \(\omega=\SI{20}{rad.s^{-1}}\). Hence,
  \[
    \lambda=\frac{2\pi}{k}=\frac{2\pi}{4}=\SI{1.57}{m},
    \qquad
    v=\frac{\omega}{k}=\frac{20}{4}=\SI{5.0}{m.s^{-1}}.
  \]
  Thus, option \textbf{C}.\end{workedsolution}
\begin{workedsolution}{12}{D}\textbf{Topic: Waves.}\par A travelling wave of the form \(kx+\omega t\) moves in the negative \(x\)-direction. Its speed is
  \[
    v=\frac{\omega}{k}=\frac{24}{6}=\SI{4}{m.s^{-1}}.
  \]
  Therefore, option \textbf{D}.\end{workedsolution}
\begin{workedsolution}{13}{B}\textbf{Topic: Waves.}\par Adjacent nodes of a standing wave are separated by half a wavelength:
  \[
    \Delta x_{\text{node-node}}=\frac{\lambda}{2}
    =\frac{0.80}{2}=\SI{0.40}{m}.
  \]
  Hence option \textbf{B}.\end{workedsolution}
\begin{workedsolution}{14}{E}\textbf{Topic: Waves.}\par For a string fixed at both ends,
  \[
    f_n=\frac{nv}{2L}.
  \]
  With \(n=3\),
  \[
    f_3=\frac{3(120)}{2(0.80)}=\frac{360}{1.60}=\SI{225}{Hz}.
  \]
  Thus, option \textbf{E}.\end{workedsolution}
\begin{workedsolution}{15}{C}\textbf{Topic: Waves.}\par The sound level is
  \[
    \beta=10\log_{10}\left(\frac{I}{I_0}\right).
  \]
  Here, \(I/I_0=10^{-6}/10^{-12}=10^6\), so
  \[
    \beta=10\log_{10}(10^6)=10(6)=\SI{60}{dB}.
  \]
  Therefore, option \textbf{C}.\end{workedsolution}
\begin{workedsolution}{16}{A}\textbf{Topic: Optics.}\par The speed is
  \[
    v=\frac{c}{n}=\frac{3.00\times10^8}{1.50}=\SI{2.00e8}{m.s^{-1}}.
  \]
  The frequency is fixed by the source and does not change at the boundary. Hence,
  \[
    \lambda=\frac{\lambda_0}{n}=\frac{600}{1.50}=\SI{400}{nm}.
  \]
  Thus, option \textbf{A}.\end{workedsolution}
\begin{workedsolution}{17}{D}\textbf{Topic: Optics.}\par Snell's law gives
  \[
    n_1\sin\theta_1=n_2\sin\theta_2.
  \]
  With \(n_1=1.00\), \(n_2=1.50\), and \(\theta_1=30.0^\circ\),
  \[
    \sin\theta_2=\frac{\sin30^\circ}{1.50}=\frac{0.5}{1.5}=\frac13.
  \]
  Therefore, \(\theta_2=\sin^{-1}(1/3)=19.5^\circ\), so option \textbf{D}.\end{workedsolution}
\begin{workedsolution}{18}{B}\textbf{Topic: Optics.}\par For small angles, the double-slit fringe spacing is
  \[
    \Delta y=\frac{\lambda L}{d}.
  \]
  Thus,
  \[
    \Delta y=\frac{(600\times10^{-9})(2.0)}{0.40\times10^{-3}}
    =3.00\times10^{-3}\ \mathrm{m}
    =\SI{3.00}{mm}.
  \]
  Hence option \textbf{B}.\end{workedsolution}
\begin{workedsolution}{19}{E}\textbf{Topic: Optics.}\par For the first single-slit minimum,
  \[
    a\sin\theta=\lambda.
  \]
  Therefore,
  \[
    \sin\theta=\frac{500\times10^{-9}}{0.20\times10^{-3}}=2.50\times10^{-3}.
  \]
  Thus, \(\theta\approx2.50\times10^{-3}\ \mathrm{rad}=0.143^\circ\). Option \textbf{E}.\end{workedsolution}
\begin{workedsolution}{20}{C}\textbf{Topic: Optics.}\par Rayleigh's criterion for a circular aperture is
  \[
    \theta_R=1.22\frac{\lambda}{D}.
  \]
  Hence,
  \[
    \theta_R=1.22\frac{550\times10^{-9}}{0.10}
    =6.71\times10^{-6}\ \mathrm{rad}.
  \]
  Therefore, option \textbf{C}.\end{workedsolution}
\subsection{Level II:\ Applied Problems}
\begin{workedsolution}{21}{B}\textbf{Topic: Oscillations.}\par At \(t=0\),
  \[
    x_0=A\cos\phi,\qquad v_0=-A\omega\sin\phi.
  \]
  Therefore,
  \[
    A^2=x_0^2+{\left(\frac{v_0}{\omega}\right)}^2
    ={(0.030)}^2+{(0.040)}^2={(0.050)}^2,
  \]
  so \(A=\SI{5.0}{cm}\). Also,
  \[
    \cos\phi=0.030/0.050=0.6,
    \qquad
    \sin\phi=-v_0/(A\omega)=0.8.
  \]
  Both are positive, so \(\phi=53.1^\circ\). Hence option \textbf{B}.\end{workedsolution}
\begin{workedsolution}{22}{D}\textbf{Topic: Oscillations.}\par For a physical pendulum,
  \[
    T=2\pi\sqrt{\frac{I}{mgh}}.
  \]
  For a uniform rod pivoted at one end, \(I=mL^2/3\) and the centre of mass is at \(h=L/2\). Thus,
  \[
    T=2\pi\sqrt{\frac{mL^2/3}{mg(L/2)}}
    =2\pi\sqrt{\frac{2L}{3g}}.
  \]
  With \(L=\SI{1.20}{m}\),
  \[
    T=2\pi\sqrt{\frac{2(1.20)}{3(9.80)}}=\SI{1.80}{s}.
  \]
  Therefore, option \textbf{D}.\end{workedsolution}
\begin{workedsolution}{23}{C}\textbf{Topic: Oscillations.}\par For three identical springs in series,
  \[
    \frac{1}{k_{\mathrm{eq}}}=\frac{1}{k}+\frac{1}{k}+\frac{1}{k}=\frac{3}{k},
  \]
  so \(k_{\mathrm{eq}}=k/3=\SI{40}{N.m^{-1}}\). Therefore,
  \[
    T=2\pi\sqrt{\frac{m}{k_{\mathrm{eq}}}}
    =2\pi\sqrt{\frac{0.36}{40}}
    =\SI{0.596}{s}.
  \]
  Hence option \textbf{C}.\end{workedsolution}
\begin{workedsolution}{24}{A}\textbf{Topic: Oscillations.}\par During the short collision, horizontal momentum is conserved:
  \[
    m_{b}v_{b}=(M+m_{b})V.
  \]
  Thus,
  \[
    V=\frac{(0.0100)(400)}{1.99+0.0100}=\SI{2.00}{m.s^{-1}}.
  \]
  After the collision, the total mass is \SI{2.00}{kg}, so
  \[
    \omega=\sqrt{\frac{k}{M+m_b}}=\sqrt{\frac{800}{2.00}}=\SI{20.0}{rad.s^{-1}}.
  \]
  Starting at equilibrium with speed \(V\), the amplitude is \(A=V/\omega=2.00/20.0=\SI{0.100}{m}\). Option \textbf{A}.\end{workedsolution}
\begin{workedsolution}{25}{E}\textbf{Topic: Oscillations.}\par Conservation of energy gives
  \[
    \frac12\kappa\theta_{\max}^2
    =\frac12\kappa\theta^2+\frac12I\dot\theta^2.
  \]
  Hence,
  \[
    \dot\theta=\sqrt{\frac{\kappa}{I}\left(\theta_{\max}^2-\theta^2\right)}.
  \]
  Substituting,
  \[
    \dot\theta=\sqrt{\frac{0.80}{0.020}\left(0.25^2-0.15^2\right)}
    =\sqrt{40(0.0400)}
    =\SI{1.26}{rad.s^{-1}}.
  \]
  Thus, option \textbf{E}.\end{workedsolution}
\begin{workedsolution}{26}{B}\textbf{Topic: Normal Modes.}\par The symmetric-system frequencies are
  \[
    \omega_1=\sqrt{\frac{k}{m}},\qquad
    \omega_2=\sqrt{\frac{k+2k_c}{m}}.
  \]
  Thus,
  \[
    \omega_1=\sqrt{\frac{100}{0.50}}=\sqrt{200}=\SI{14.1}{rad.s^{-1}},
  \]
  \[
    \omega_2=\sqrt{\frac{100+100}{0.50}}=\sqrt{400}=\SI{20.0}{rad.s^{-1}}.
  \]
  Therefore, option \textbf{B}.\end{workedsolution}
\begin{workedsolution}{27}{D}\textbf{Topic: Normal Modes.}\par From the in-phase frequency,
  \[
    k=m\omega_1^2=2.0{(10)}^2=\SI{200}{N.m^{-1}}.
  \]
  From the out-of-phase frequency,
  \[
    k+2k_c=m\omega_2^2=2.0{(14)}^2=\SI{392}{N.m^{-1}}.
  \]
  Therefore,
  \[
    2k_c=392-200=192,
    \qquad k_c=\SI{96}{N.m^{-1}}.
  \]
  Hence option \textbf{D}.\end{workedsolution}
\begin{workedsolution}{28}{C}\textbf{Topic: Normal Modes.}\par Equating components gives
  \[
    C_1+C_2=A,
    \qquad
    C_1-C_2=0.
  \]
  The second equation gives \(C_1=C_2\). Substituting into the first gives \(2C_1=A\), so
  \[
    C_1=C_2=\frac{A}{2}.
  \]
  Therefore, option \textbf{C}.\end{workedsolution}
\begin{workedsolution}{29}{A}\textbf{Topic: Normal Modes.}\par Direct substitution gives
  \[
    q_1=x_1+x_2=3.0+(-1.0)=\SI{2.0}{cm},
  \]
  \[
    q_2=x_1-x_2=3.0-(-1.0)=\SI{4.0}{cm}.
  \]
  Thus, option \textbf{A}. The positive \(q_2\) indicates a substantial out-of-phase component.\end{workedsolution}
\begin{workedsolution}{30}{E}\textbf{Topic: Normal Modes.}\par The normal-mode condition is
  \[
    \det(K-m\omega^2I)=0,
  \]
  with
  \[
    K-m\omega^2I=
    \begin{bmatrix}
      k+k_c-m\omega^2 & -k_c            \\
      -k_c            & k+k_c-m\omega^2
    \end{bmatrix}.
  \]
  Its determinant is
  \[
    {(k+k_c-m\omega^2)}^2-(-k_c)(-k_c)
    ={(k+k_c-m\omega^2)}^2-k_c^2.
  \]
  Hence option \textbf{E}.\end{workedsolution}
\begin{workedsolution}{31}{B}\textbf{Topic: Waves.}\par Here \(k=\SI{5}{rad.m^{-1}}\) and \(\omega=\SI{100}{rad.s^{-1}}\), so
  \[
    v=\frac{\omega}{k}=\SI{20}{m.s^{-1}}.
  \]
  Since \(v=\sqrt{T/\mu}\),
  \[
    T=\mu v^2=0.010{(20)}^2=\SI{4.0}{N}.
  \]
  The average power is
  \[
    P_{\mathrm{av}}=\frac12\mu\omega^2A^2v
    =\frac12{(0.010)}{(100)}^2{(0.004)}^2{(20)}
    =\SI{0.016}{W}.
  \]
  Thus, option \textbf{B}.\end{workedsolution}
\begin{workedsolution}{32}{C}\textbf{Topic: Waves.}\par For a closed-open pipe,
  \[
    f_m=\frac{(2m+1)v}{4L},\qquad m=0,1,2,\ldots
  \]
  The fundamental is
  \[
    f_1=\frac{340}{4(0.85)}=\SI{100}{Hz}.
  \]
  Only odd harmonics occur, so the next two are \(3f_1=\SI{300}{Hz}\) and \(5f_1=\SI{500}{Hz}\). Therefore, option \textbf{C}.\end{workedsolution}
\begin{workedsolution}{33}{D}\textbf{Topic: Waves.}\par The increased tension raises the frequency to \(505\ \mathrm{Hz}\). For the same string length and linear density,
  \[
    f\propto\sqrt{T}.
  \]
  Therefore,
  \[
    \frac{T_f}{T_i}={\left(\frac{f_f}{f_i}\right)}^2
    ={\left(\frac{505}{500}\right)}^2=1.0201.
  \]
  The fractional increase is \(1.0201-1=0.0201\). Hence option \textbf{D}.\end{workedsolution}
\begin{workedsolution}{34}{A}\textbf{Topic: Waves.}\par For a moving source approaching a stationary observer,
  \[
    f'=f\frac{v}{v-v_s}.
  \]
  Thus,
  \[
    f'=600\frac{343}{343-30}
    =600\frac{343}{313}
    =\SI{658}{Hz}\quad\text{(approximately)}.
  \]
  The observed frequency is higher because the source approaches. Option \textbf{A}.\end{workedsolution}
\begin{workedsolution}{35}{E}\textbf{Topic: Waves.}\par The Mach angle satisfies
  \[
    \sin\mu=\frac{1}{M}.
  \]
  For \(M=1.50\),
  \[
    \mu=\sin^{-1}\left(\frac{1}{1.50}\right)
    =\sin^{-1}(0.6667)
    =41.8^\circ.
  \]
  Therefore, option \textbf{E}.\end{workedsolution}
\begin{workedsolution}{36}{B}\textbf{Topic: Optics.}\par The sheet introduces additional optical path
  \[
    \Delta L=(n-1)t.
  \]
  The number of fringe spacings shifted is
  \[
    N=\frac{\Delta L}{\lambda}
    =\frac{(1.60-1)(5.0\times10^{-6})}{600\times10^{-9}}
    =\frac{3.0\times10^{-6}}{0.60\times10^{-6}}=5.
  \]
  Hence option \textbf{B}.\end{workedsolution}
\begin{workedsolution}{37}{C}\textbf{Topic: Optics.}\par An interference maximum is missing when it coincides with a single-slit minimum. The condition is
  \[
    m=p\frac{d}{a}.
  \]
  Since \(d/a=4\),
  \[
    m=4p=4,8,12,\ldots
  \]
  Therefore, option \textbf{C}.\end{workedsolution}
\begin{workedsolution}{38}{D}\textbf{Topic: Optics.}\par The line spacing is
  \[
    d=\frac{1}{500\ \mathrm{mm}^{-1}}=0.00200\ \mathrm{mm}=\SI{2.00e-6}{m}.
  \]
  For \(m=2\),
  \[
    \sin\theta_2=\frac{2\lambda}{d}
    =\frac{2(600\times10^{-9})}{2.00\times10^{-6}}=0.600,
  \]
  so \(\theta_2=36.9^\circ\). Also,
  \[
    m_{\max}=\left\lfloor\frac{d}{\lambda}\right\rfloor
    =\lfloor3.33\rfloor=3.
  \]
  Hence option \textbf{D}.\end{workedsolution}
\begin{workedsolution}{39}{A}\textbf{Topic: Optics.}\par For critical incidence from glass to air,
  \[
    \sin\theta_c=\frac{n_{\text{air}}}{n_{\text{glass}}}=\frac{1}{1.50},
  \]
  so \(\theta_c=41.8^\circ\). For Brewster incidence from air to glass,
  \[
    \tan\theta_B=\frac{n_{\text{glass}}}{n_{\text{air}}}=1.50,
  \]
  so \(\theta_B=56.3^\circ\). Therefore, option \textbf{A}.\end{workedsolution}
\begin{workedsolution}{40}{E}\textbf{Topic: Optics.}\par Bragg's law uses \(\theta\), not the displayed angle \(2\theta\). Thus \(\theta=25.0^\circ\). For first order,
  \[
    \lambda=2d\sin\theta.
  \]
  Therefore,
  \[
    d=\frac{0.154}{2\sin25.0^\circ}\ \mathrm{nm}
    =\SI{0.182}{nm}.
  \]
  Hence option \textbf{E}.\end{workedsolution}
\subsection{Level III:\ Multi-step Challenges}
\begin{workedsolution}{41}{C}\textbf{Topic: Oscillations.}\par Using \(x=A\cos(\omega t+\phi)\) at \(t=0\),
  \[
    \cos\phi=\frac{x_0}{A}=\frac{0.030}{0.060}=\frac12.
  \]
  Thus \(\phi=\pm\pi/3\). The velocity is
  \[
    v=-A\omega\sin(\omega t+\phi).
  \]
  For \(v(0)>0\), we require \(\sin\phi<0\), so \(\phi=-\pi/3\). Therefore,
  \[
    x=0.060\cos(8t-\pi/3),
  \]
  which is option \textbf{C}.\end{workedsolution}
\begin{workedsolution}{42}{B}\textbf{Topic: Oscillations.}\par The SHM speed-position relation is
  \[
    v=\omega\sqrt{A^2-x^2}.
  \]
  Here,
  \[
    \sqrt{A^2-x^2}=\sqrt{{(0.050)}^2-{(0.030)}^2}
    =\sqrt{0.0016}=\SI{0.040}{m}.
  \]
  Thus,
  \[
    \omega=\frac{0.40}{0.040}=\SI{10}{rad.s^{-1}},
    \qquad
    T=\frac{2\pi}{\omega}=\SI{0.628}{s}.
  \]
  Therefore, option \textbf{B}.\end{workedsolution}
\begin{workedsolution}{43}{D}\textbf{Topic: Oscillations.}\par The centre of mass is at \(L/2\) from the end, while the pivot is at \(L/4\), so \(h=L/4\). The moment of inertia about the pivot is found by the parallel-axis theorem:
  \[
    I=I_{\mathrm{cm}}+mh^2
    =\frac{mL^2}{12}+m{\left\{\frac{L}{4}\right\}}^2
    =\frac{7mL^2}{48}.
  \]
  Therefore,
  \[
    T=2\pi\sqrt{\frac{I}{mgh}}
    =2\pi\sqrt{\frac{7L}{12g}}
    =\SI{1.68}{s}.
  \]
  Hence option \textbf{D}.\end{workedsolution}
\begin{workedsolution}{44}{A}\textbf{Topic: Oscillations.}\par Write \(x=A\cos(\omega t+\phi)\). At \(t=0\), \(x=A/2\), so \(\cos\phi=1/2\). Moving toward equilibrium from the positive side means \(v<0\), hence \(\sin\phi>0\) and \(\phi=\pi/3\). The opposite turning point occurs when the phase first reaches \(\pi\):
  \[
    \omega t+\frac{\pi}{3}=\pi
    \quad\Rightarrow\quad
    \omega t=\frac{2\pi}{3}.
  \]
  Since \(T=2\pi/\omega\), \(t=T/3\). Thus, option \textbf{A}.\end{workedsolution}
\begin{workedsolution}{45}{E}\textbf{Topic: Oscillations.}\par During sticking, horizontal momentum is conserved because the clay initially has no horizontal momentum:
  \[
    Mv_0=(M+m)V.
  \]
  Thus,
  \[
    V=\frac{(1.00)(2.00)}{1.25}=\SI{1.60}{m.s^{-1}}.
  \]
  The new angular frequency is
  \[
    \omega'=\sqrt{\frac{k}{M+m}}=\sqrt{\frac{100}{1.25}}=\SI{8.94}{rad.s^{-1}}.
  \]
  At equilibrium all post-collision oscillator energy is kinetic, so \(A'=V/\omega'=1.60/8.94=\SI{0.179}{m}\). Option \textbf{E}.\end{workedsolution}
\begin{workedsolution}{46}{C}\textbf{Topic: Normal Modes.}\par The frequencies are
  \[
    \omega_1=\sqrt{\frac{k}{m}}=2,
    \qquad
    \omega_2=\sqrt{\frac{k+2k_c}{m}}=\sqrt{16}=4\ \mathrm{rad\,s^{-1}}.
  \]
  The initial displacement \({[A,0]}^T\) contains equal modal coefficients:
  \[
    x_1=\frac{A}{2}(\cos\omega_1t+\cos\omega_2t),
    \quad
    x_2=\frac{A}{2}(\cos\omega_1t-\cos\omega_2t).
  \]
  At \(t=\pi/2\), \(\cos(2t)=\cos\pi=-1\) and \(\cos(4t)=\cos2\pi=1\). Hence \(x_1=0\) and \(x_2=-A\). Option \textbf{C}.\end{workedsolution}
\begin{workedsolution}{47}{B}\textbf{Topic: Normal Modes.}\par First convert to angular frequencies:
  \[
    \omega_1=2\pi f_1=4\pi,
    \qquad
    \omega_2=2\pi f_2=6\pi.
  \]
  Then
  \[
    k=m\omega_1^2=0.50{(4\pi)}^2=\SI{79.0}{N.m^{-1}}.
  \]
  Also,
  \[
    k+2k_c=m\omega_2^2,
  \]
  so
  \[
    k_c=\frac{m(\omega_2^2-\omega_1^2)}{2}
    =\frac{0.50[{(6\pi)}^2-{(4\pi)}^2]}{2}
    =\SI{49.3}{N.m^{-1}}.
  \]
  Therefore, option \textbf{B}.\end{workedsolution}
\begin{workedsolution}{48}{D}\textbf{Topic: Normal Modes.}\par For the initial condition \({[A,0]}^T\), the displacement of the first mass can be written
  \[
    x_1=A\cos\left(\frac{\omega_1+\omega_2}{2}t\right)
    \cos\left(\frac{\omega_2-\omega_1}{2}t\right).
  \]
  Its amplitude envelope first becomes zero when
  \[
    \cos\left(\frac{\Delta\omega}{2}t\right)=0.
  \]
  The first solution is \((\Delta\omega/2)t=\pi/2\), hence
  \[
    t=\frac{\pi}{\Delta\omega}=\frac{\pi}{12-10}=\SI{1.57}{s}.
  \]
  At this time the amplitude has transferred to the other mass. Option \textbf{D}.\end{workedsolution}
\begin{workedsolution}{49}{A}\textbf{Topic: Normal Modes.}\par Evaluate the energy at a turning point, where kinetic energy is zero. In phase, each outer spring stores \(\tfrac12kA^2\), while the coupling spring is unstrained:
  \[
    E_{\mathrm{in}}=kA^2.
  \]
  Out of phase, the outer springs still store \(kA^2\) in total, but the coupling spring changes length by \(2A\):
  \[
    E_{\mathrm{out}}=kA^2+\frac12k_c{(2A)}^2={(k+2k_c)}A^2.
  \]
  Thus,
  \[
    \frac{E_{\mathrm{out}}}{E_{\mathrm{in}}}=\frac{k+2k_c}{k}=\frac{100+100}{100}=2.
  \]
  Hence option \textbf{A}.\end{workedsolution}
\begin{workedsolution}{50}{E}\textbf{Topic: Normal Modes.}\par Acting with \(K\) on the mode vectors gives
  \[
    K\begin{bmatrix}1\\1\end{bmatrix}=4\begin{bmatrix}1\\1\end{bmatrix},
    \qquad
    K\begin{bmatrix}1\\-1\end{bmatrix}=8\begin{bmatrix}1\\-1\end{bmatrix}.
  \]
  For unit masses, \(\omega_1=\sqrt4=2\) and \(\omega_2=\sqrt8=2\sqrt2\). For the displacement decomposition,
  \[
    C_1+C_2=3,\qquad C_1-C_2=1,
  \]
  so \(C_1=2\ \mathrm{cm}\) and \(C_2=1\ \mathrm{cm}\). Thus, option \textbf{E}.\end{workedsolution}
\begin{workedsolution}{51}{C}\textbf{Topic: Waves.}\par Represent the three contributions as phasors:
  \[
    A_R=A\left(1+e^{i\pi/3}+e^{i2\pi/3}\right).
  \]
  Using \(e^{i\pi/3}=\tfrac12+i\tfrac{\sqrt3}{2}\) and \(e^{i2\pi/3}=-\tfrac12+i\tfrac{\sqrt3}{2}\),
  \[
    1+e^{i\pi/3}+e^{i2\pi/3}=1+i\sqrt3=2e^{i\pi/3}.
  \]
  Therefore, the resultant amplitude is \(2A=\SI{10.0}{mm}\) and its phase is \(\pi/3\). Option \textbf{C}.\end{workedsolution}
\begin{workedsolution}{52}{B}\textbf{Topic: Waves.}\par The total time is the fall time plus the upward sound-travel time:
  \[
    \sqrt{\frac{2h}{g}}+\frac{h}{v}=4.00.
  \]
  Let \(u=\sqrt{h}\). Then
  \[
    \frac{u^2}{343}+\sqrt{\frac{2}{9.80}}u-4.00=0.
  \]
  The positive root is \(u=8.399\), so
  \[
    h=u^2=\SI{70.5}{m}.
  \]
  As a check, the fall time is \SI{3.794}{s} and the sound time is \SI{0.206}{s}, summing to \SI{4.00}{s}. Hence option \textbf{B}.\end{workedsolution}
\begin{workedsolution}{53}{E}\textbf{Topic: Waves.}\par There are two Doppler shifts. First, the moving car is an observer approaching the stationary source:
  \[
    f_1=f_0\frac{v+u}{v}.
  \]
  The car then acts as a moving source approaching the detector:
  \[
    f_2=f_1\frac{v}{v-u}.
  \]
  Therefore,
  \[
    f_2=f_0\frac{v+u}{v-u}
    =40.0\frac{343+30}{343-30}\ \mathrm{kHz}
    =\SI{47.7}{kHz}.
  \]
  Thus, option \textbf{E}.\end{workedsolution}
\begin{workedsolution}{54}{A}\textbf{Topic: Waves.}\par The local amplitude is the magnitude of the position factor:
  \[
    A_{\mathrm{local}}=0.080\left|\sin\left(\frac{\pi(0.50)}{2}\right)\right|
    =0.080\sin\frac{\pi}{4}
    =\SI{0.0566}{m}.
  \]
  At \(t=0.0625=1/16\ \mathrm{s}\),
  \[
    \cos(4\pi t)=\cos\frac{\pi}{4}=\frac{\sqrt2}{2}.
  \]
  Thus,
  \[
    y=0.080\left(\frac{\sqrt2}{2}\right)\left(\frac{\sqrt2}{2}\right)
    =\SI{0.0400}{m}.
  \]
  Hence option \textbf{A}.\end{workedsolution}
\begin{workedsolution}{55}{D}\textbf{Topic: Waves.}\par Far from the sources, the path difference is \(\Delta L=d\cos\theta\) (an equivalent sine convention gives the same count). Maxima require
  \[
    d\cos\theta=m\lambda.
  \]
  Because \(d/\lambda=1.50/0.400=3.75\), the allowed integers are
  \[
    m=-3,-2,-1,0,1,2,3,
  \]
  seven values. For each value, \(|m\lambda/d|<1\), so \(\cos\theta=m\lambda/d\) has two positions during a full circuit. Therefore, the total is \(7\times2=14\). Option \textbf{D}.\end{workedsolution}
\begin{workedsolution}{56}{B}\textbf{Topic: Optics.}\par Substitute the specified values:
  \[
    {\left(\frac{\sin\beta}{\beta}\right)}^2
    ={\left(\frac{1}{\pi/2}\right)}^2
    =\frac{4}{\pi^2},
  \]
  and
  \[
    \cos^2\alpha=\cos^2\frac{\pi}{3}={\left(\frac12\right)}^2=\frac14.
  \]
  Therefore,
  \[
    \frac{I}{I_0}=\frac{4}{\pi^2}\cdot\frac14=\frac{1}{\pi^2}\approx0.101.
  \]
  Hence option \textbf{B}.\end{workedsolution}
\begin{workedsolution}{57}{C}\textbf{Topic: Optics.}\par The required resolving power is
  \[
    R=\frac{\lambda}{\Delta\lambda}
    =\frac{600.0}{0.3}=2000.
  \]
  For a grating, \(R=mN\). In second order,
  \[
    N=\frac{R}{m}=\frac{2000}{2}=1000.
  \]
  At \(500\) lines per millimetre, the illuminated width is
  \[
    W=\frac{1000}{500\ \mathrm{mm}^{-1}}=\SI{2.00}{mm}.
  \]
  Therefore, option \textbf{C}.\end{workedsolution}
\begin{workedsolution}{58}{D}\textbf{Topic: Optics.}\par The parameter is
  \[
    \beta=\frac{\pi a\sin\theta}{\lambda}.
  \]
  Thus,
  \[
    \sin\theta=\frac{\beta\lambda}{\pi a}
    =\frac{4.493(600\times10^{-9})}{\pi(0.20\times10^{-3})}
    =4.29\times10^{-3}.
  \]
  Hence \(\theta\approx0.246^\circ\). The first secondary maximum has the standard intensity \(I\approx0.047I_{\max}\), or \(4.7\%\). Therefore, option \textbf{D}.\end{workedsolution}
\begin{workedsolution}{59}{A}\textbf{Topic: Optics.}\par Their small angular separation is
  \[
    \theta\approx\frac{s}{L}=\frac{1.50}{100\,000}=1.50\times10^{-5}\ \mathrm{rad}.
  \]
  Rayleigh's criterion requires
  \[
    \theta=1.22\frac{\lambda}{D}.
  \]
  Therefore,
  \[
    D=1.22\frac{550\times10^{-9}}{1.50\times10^{-5}}
    =4.47\times10^{-2}\ \mathrm{m}
    =\SI{4.47}{cm}.
  \]
  Hence option \textbf{A}.\end{workedsolution}
\begin{workedsolution}{60}{E}\textbf{Topic: Optics.}\par The spacing is
  \[
    d=\frac{1}{300\ \mathrm{mm}^{-1}}=3.333\times10^{-6}\ \mathrm{m}.
  \]
  The largest possible positive order is
  \[
    m_{\max}=\left\lfloor\frac{d}{\lambda}\right\rfloor
    =\left\lfloor\frac{3.333\times10^{-6}}{500\times10^{-9}}\right\rfloor
    =\lfloor6.667\rfloor=6.
  \]
  Because \(a=d/3\), we have \(d/a=3\), so missing orders satisfy \(m=3p\): \(m=3\) and \(m=6\) within the allowed range. The observable positive orders are therefore \(1,2,4,5\), four on each side. Including the central order gives \(2(4)+1=9\). Hence option \textbf{E}.\end{workedsolution}
\clearpage
\section{Compact Formula Checklist}
\begin{multicols}{2}
  \subsection*{Oscillations}
  \[
    F=-kx,\quad \omega=\sqrt{\frac{k}{m}},\quad T=\frac{2\pi}{\omega},
  \]
  \[
    x=A\cos(\omega t+\phi),\quad v=-A\omega\sin(\omega t+\phi),
  \]
  \[
    a=-\omega^2x,\quad E=\frac12kA^2,\quad v^2=\omega^2(A^2-x^2).
  \]
  \[
    T_{\text{simple pendulum}}=2\pi\sqrt{\frac{L}{g}},\qquad
    T_{\text{physical}}=2\pi\sqrt{\frac{I}{mgh}}.
  \]

  \subsection*{Normal Modes}
  \[
    M\ddot{\mathbf{x}}+K\mathbf{x}=0,
    \quad (K-\omega^2M)\mathbf{A}=0,
    \quad \det(K-\omega^2M)=0.
  \]
  For the symmetric two-mass model,
  \[
    \omega_1=\sqrt{\frac{k}{m}},\quad \mathbf{A}_1\propto\begin{bmatrix}1\\1\end{bmatrix},
  \]
  \[
    \omega_2=\sqrt{\frac{k+2k_c}{m}},\quad \mathbf{A}_2\propto\begin{bmatrix}1\\-1\end{bmatrix}.
  \]

  \subsection*{Waves and Sound}
  \[
    k=\frac{2\pi}{\lambda},\quad \omega=2\pi f,\quad v=f\lambda=\frac{\omega}{k}.
  \]
  \[
    v_{\text{string}}=\sqrt{\frac{T}{\mu}},\quad P_{\mathrm{av}}=\frac12\mu\omega^2A^2v.
  \]
  \[
    f_n=\frac{nv}{2L}\quad \text{(fixed-fixed/open-open)},
  \]
  \[
    f_m=\frac{(2m+1)v}{4L}\quad \text{(closed-open)}.
  \]
  \[
    f'=f\left(\frac{v\pm v_o}{v\mp v_s}\right),\quad
    \sin\mu=\frac{1}{M}.
  \]

  \subsection*{Optics}
  \[
    n_1\sin\theta_1=n_2\sin\theta_2,\quad v=\frac{c}{n},\quad \lambda=\frac{\lambda_0}{n}.
  \]
  \[
    \Delta y_{\text{double slit}}=\frac{\lambda L}{d},\quad
    \text{single-slit minima: }a\sin\theta=m\lambda.
  \]
  \[
    \theta_R=1.22\frac{\lambda}{D},\quad
    \text{grating maxima: }d\sin\theta=m\lambda.
  \]
  \[
    R=mN,\quad n\lambda=2d\sin\theta\ \text{(Bragg)}.
  \]
\end{multicols}

\end{document}
